\section{Arbitrary-factor periodic rigidity}
\label{sec:factor-rigidity}

Source proximality becomes useful only after its quantifiers survive the
factor map.  We therefore separate factor permanence from periodic-orbit
separation.  This also makes clear that the target need not be symbolic.

\subsection{Proximality passes to every lawful factor}

Fix a compact metrizable space \(Y\), a homeomorphism \(S:Y\to Y\), and a
continuous surjection \(\pi:\Xsf\twoheadrightarrow Y\) satisfying
\begin{equation}
  \pi\circ\sigma^n=S^n\circ\pi
  \qquad(n\in\mathbb Z).
  \label{eq:factor-equivariance}
\end{equation}

\begin{lemma}[Factor permanence]
\label{lem:factor-proximal}
The target system \((Y,S)\) is proximal.
\end{lemma}

\begin{proof}
Let \(y_1,y_2\in Y\).  Surjectivity supplies lifts
\(x_1,x_2\in\Xsf\) with \(\pi(x_i)=y_i\).  Because \(\Xsf\) is compact,
continuity of \(\pi\) is uniform: for every \(\varepsilon>0\), some
\(\delta>0\) satisfies
\[
  d_X(u,v)<\delta
  \Longrightarrow
  d_Y(\pi u,\pi v)<\varepsilon.
\]
By \cref{cor:source-proximal}, choose \(n\geq0\) with
\(d_X(\sigma^n x_1,\sigma^n x_2)<\delta\).  Equivariance gives
\[
  d_Y(S^n y_1,S^n y_2)
  =d_Y(\pi\sigma^n x_1,\pi\sigma^n x_2)<\varepsilon.
\]
As \(\varepsilon\) is arbitrary, \((y_1,y_2)\) is proximal.
\end{proof}

Each hypothesis has a distinct role.  Surjectivity lifts arbitrary target
pairs.  Continuity and source compactness yield uniform control.  Equivariance
compares the same time in source and target.  The theorem does not extend by
wordplay to a map missing any of these properties.

\subsection{Finite-orbit separation}

Equivariance anchors the fixed point
\begin{equation}
  y_0:=\pi(0^{\mathbb Z}),
  \qquad Sy_0=y_0.
  \label{eq:fixed-anchor}
\end{equation}

\begin{lemma}[Periodic separation]
\label{lem:periodic-separation}
Let \((Y,S)\) be a proximal metrizable \(\mathbb Z\)-system with a fixed
point \(y_0\).  Then \(y_0\) is the only periodic point.
\end{lemma}

\begin{proof}
Suppose \(y\) has least period \(r\).  If \(r=1\) and \(y\neq y_0\), then
\[
  d_Y(S^n y,S^n y_0)=d_Y(y,y_0)>0
  \qquad(n\geq0),
\]
so the pair \((y,y_0)\) is not proximal.

Now suppose \(r>1\).  The points \(S^k y\) and \(S^{k+1}y\) are distinct
for every \(0\leq k<r\).  The finite minimum
\begin{equation}
  \delta_r:=\min_{0\leq k<r}
  d_Y(S^k y,S^{k+1}y)
  \label{eq:periodic-gap}
\end{equation}
is therefore positive.  Reducing \(n\) modulo \(r\) shows
\[
  d_Y(S^n y,S^n(Sy))\geq\delta_r
  \qquad(n\geq0).
\]
Thus \((y,Sy)\) is not proximal.  Both cases contradict proximality unless
\(y=y_0\).
\end{proof}

\begin{figure}[t]
\centering
\begin{tikzpicture}[
  node distance=10mm and 15mm,
  system/.style={draw,rounded corners=2pt,align=center,minimum height=14mm,
    text width=0.31\textwidth,inner sep=5pt},
  case/.style={draw,rounded corners=2pt,align=center,text width=0.36\textwidth,
    inner sep=5pt,font=\small},
  arrow/.style={-{Latex[length=2.5mm]},thick}
]
\node[system,draw=sourceblue,fill=sourceblue!7] (source)
  {\(\Xsf\)\\compact, two-sided, proximal};
\node[system,draw=factorblue,fill=factorblue!7,right=of source] (target)
  {arbitrary compact metrizable \(Y\)\\\(S\) a homeomorphism};
\draw[arrow,factorblue] (source) -- (target);
\node[font=\small,align=center,fill=white,inner sep=2pt,
  text width=0.28\textwidth]
  at ($(source.north)!0.5!(target.north)+(0,9mm)$)
  {continuous, onto,\\\(\mathbb Z\)-equivariant \(\pi\)};

\node[case,draw=stopred,fill=stopred!5,below=13mm of source]
  (fixedcase) {\textbf{Second fixed point:}
  \(d(S^n y,S^n y_0)=d(y,y_0)>0\).};
\node[case,draw=stopred,fill=stopred!5,right=of fixedcase]
  (cyclecase) {\textbf{Least period \(r>1\):}
  \(d(S^n y,S^n Sy)\geq\delta_r>0\).};
\draw[arrow,stopred] (target.south) -- ++(0,-4mm) -| (fixedcase.north);
\draw[arrow,stopred] (target.south) -- ++(0,-4mm) -| (cyclecase.north);

\coordinate (midcases) at ($(fixedcase.south)!0.5!(cyclecase.south)$);
\node[draw=proofgreen,fill=proofgreen!6,rounded corners=2pt,
  below=11mm of midcases,
  text width=0.55\textwidth,align=center,inner sep=5pt] (conclusion)
  {both cases contradict proximality; therefore
  \(\Per(Y,S)=\{\pi(0^{\mathbb Z})\}\)};
\draw[arrow,proofgreen] (fixedcase) -- (conclusion);
\draw[arrow,proofgreen] (cyclecase) -- (conclusion);
\end{tikzpicture}
\caption{The arbitrary-factor quantifier and the two periodic-separation
cases.  No symbolic presentation of \(Y\) is used.}
\label{fig:factor-separation}
\end{figure}


\subsection{The theorem}

\begin{theorem}[Squarefree factor periodic rigidity]
\label{thm:factor-rigidity}
Let \(\Xsf\) be the exact source in
\cref{eq:source-definition-intro}.  For every compact metrizable
\(\mathbb Z\)-system \((Y,S)\) and every continuous surjective equivariant
map \(\pi:\Xsf\twoheadrightarrow Y\),
\[
  \Per(Y,S)=\{\pi(0^{\mathbb Z})\}.
\]
\end{theorem}

\begin{proof}
By \cref{eq:fixed-anchor}, \(y_0=\pi(0^{\mathbb Z})\) is fixed.
\Cref{cor:source-proximal,lem:factor-proximal} make \((Y,S)\) proximal, and
\cref{lem:periodic-separation} excludes every periodic point other than
\(y_0\).
\end{proof}

The word ``every'' in \cref{thm:factor-rigidity} refers to the entire lawful
factor category, not to a finite list of symbolic quotients.  The proof never
examines a target alphabet or a coding radius.  Conversely, the theorem is
not a universal statement that factors of aperiodic systems are aperiodic.
It succeeds because the frozen source is proximal; deleting that property
invalidates the finite-orbit contradiction.

\subsection{Assumption ledger}

\begin{table}[t]
\centering
\small
\caption{Where the factor assumptions enter.}
\label{tab:factor-assumptions}
\begin{tabularx}{\textwidth}{@{}l>{\raggedright\arraybackslash}X
  >{\raggedright\arraybackslash}X@{}}
\toprule
Assumption & Use in the proof & What deletion permits \\
\midrule
all prime-square exclusions & fresh coprime moduli for every window & finite
  approximants with cycles \\
continuity & small source distance gives small target distance & arbitrary
  observations need not preserve proximality \\
surjectivity & lifts every target pair & a partial image says nothing about
  all target points \\
equivariance & aligns source and target time & time-varying or unrelated maps \\
compact metric source & uniform continuity & merely pointwise control is
  insufficient for the stated argument \\
homeomorphic \(\mathbb Z\)-action & exact two-sided factor category & changed
  semigroup/induced dynamics \\
\bottomrule
\end{tabularx}
\end{table}

